% =====================================================================
% Related Work (CVPR style, ~1 column, 4 subsections)
% Each subsection ends with one sentence that states the gap and how
% SearchCop addresses it. The final paragraph contrasts SearchCop with
% the closest tool-using vision-language agents.
% =====================================================================
\section{Related Work}
\label{sec:related}

\subsection{Video Re-Identification}
\label{ssec:rel_videoreid}
The dominant body of work on person retrieval treats the problem as
representation learning over single-camera tracklets. Image-level methods
such as TransReID~\cite{he2021transreid} and CLIP-ReID~\cite{li2023clipreid}
introduced strong transformer or vision-language pretraining backbones,
while video-level methods such as AP3D~\cite{gu2020ap3d},
GRL~\cite{liu2021grl}, and STMN~\cite{eom2021stmn} aggregate temporal
information through 3D convolutions, attention, or memory-based pooling.
Standard benchmarks include MARS~\cite{zheng2016mars} and
DukeMTMC-VideoReID~\cite{wu2018duke}; both fix the gallery to a single
modality and assume the system retrieves Top-$K$ in one shot.
\textbf{Gap.} These approaches optimize a single similarity function and
deliver no \emph{decision policy}: the system cannot, for example, decide
to switch modality when the query is occluded, nor explain why a Top-1
prediction was chosen. \method{} preserves these encoders as
\emph{tools} but lifts modality selection and aggregation into an explicit
LLM-driven plan.

\subsection{Multi-Camera Person Search and Spatio-Temporal Reasoning}
\label{ssec:rel_mcps}
Multi-camera person search jointly tackles detection and re-identification
across multiple cameras, with PRW~\cite{zheng2017prw} and
CUHK-SYSU~\cite{xiao2017cuhksysu} as canonical benchmarks; recent
MEVID~\cite{davila2023mevid} explicitly targets multi-environment,
long-time-span video. A handful of works exploit
\emph{spatio-temporal} priors -- camera topology, transit-time histograms,
or graph propagation -- to refine retrieval~\cite{lv2018global,wang2019stp}.
Yet, both lines treat each gallery item independently and never plan a
\emph{sequence} of retrieval actions; topology, when used at all, is
applied as a fixed post-filter rather than a tool the system can choose
to invoke. \textbf{Gap.} The decision of \emph{when} to consult topology,
\emph{which} encoder to apply, and \emph{whether} the answer is good
enough is left to manual heuristics. \method{} promotes each of these
decisions to first-class actions in an LLM agent's tool space.

\subsection{Clothing-Change Re-ID and Gait Recognition}
\label{ssec:rel_cc_gait}
When clothing changes between camera views, appearance-based features
become unreliable. Two complementary lines emerged. (i) Clothing-aware
representation learning, exemplified by CAL~\cite{gu2022cal} on
PRCC~\cite{yang2019prcc} and LTCC~\cite{qian2020ltcc}, learns features
adversarially insensitive to clothes. (ii) Gait recognition exploits
walking dynamics rather than appearance: GaitSet~\cite{chao2019gaitset},
GaitPart~\cite{fan2020gaitpart}, and the OpenGait
family~\cite{fan2023opengait} demonstrate strong cross-view performance
on silhouette inputs. Both lines, however, presume an \emph{a priori}
choice of modality at indexing time, and offline late fusion
typically falls back to fixed weighted sums.
\textbf{Gap.} Real surveillance queries rarely come with a clean modality
hint -- a frontal occluded crop calls for Re-ID, a clean back-view
tracklet calls for gait. \method{} keeps both encoders as on-demand
tools and lets the Planner decide which one to trust per query, with the
Critic invalidating the call when modality disagreement is detected.

\subsection{LLM Agents and Tool Use}
\label{ssec:rel_agents}
Tool-using LLM agents are a fast-growing line of work originating from
ReAct~\cite{yao2023react} and Toolformer~\cite{schick2023toolformer},
which interleave reasoning and acting via natural-language traces and
self-supervised tool annotation, respectively. Building on these,
HuggingGPT~\cite{shen2023hugginggpt} orchestrates Hugging Face models for
multi-modal tasks; Visual Programming~\cite{gupta2023visprog} and
Chameleon~\cite{lu2023chameleon} compose vision tools into programs for
single-image VQA and reasoning. Beyond, AgentBench~\cite{liu2024agentbench}
and ToolBench~\cite{qin2024toolbench} measure agent capability on coding,
web, and API tasks. \textbf{Gap.} Existing agent benchmarks focus on
\emph{capability} (can the agent solve a task?) rather than
\emph{retrieval faithfulness} (does the agent's reasoning match its tool
output?), and none of these systems tackle large-scale multi-camera
gallery retrieval. \method{} adapts the tool-using paradigm to person
search and additionally proposes \emph{Reasoning Faithfulness} and
\emph{Tool-call Efficiency} as new evaluation primitives that future
agent-based retrieval systems can adopt.

\subsection{Vision-Language Agents for Video and Surveillance}
\label{ssec:rel_vlmagents}
Recent video-centric multimodal LLMs such as
VideoChat~\cite{li2023videochat}, Video-LLaMA~\cite{zhang2023videollama},
Video-ChatGPT~\cite{maaz2023videochatgpt}, and
TimeChat~\cite{ren2024timechat} push LLMs to understand long videos via
captioning or temporal QA. Closer to retrieval, AssistGPT~\cite{gao2023assistgpt}
plans modality-specific calls for video QA, and a small literature on
LLM-driven surveillance studies anomaly description. \textbf{Gap.} These
systems address \emph{describing} or \emph{answering questions about} a
single video, not \emph{searching} a multi-camera gallery for the
\emph{same person} under spatial-temporal constraints; they also have no
notion of self-correcting retrieval. To our knowledge, \method{} is the
first LLM agent that targets this exact problem.

\paragraph{Positioning.} Table~\ref{tab:positioning} contrasts \method{}
with the most relevant prior work along five axes: multi-camera support,
gait+Re-ID modality fusion, explicit decision making, interpretable
output, and evaluation on MEVID. \method{} is the first system to be
positive on all five.

\begin{table*}[t]
\centering
\small
\caption{Positioning of \method{} against representative prior work.
\textbullet\ = supported, $\circ$ = partial, blank = unsupported.}
\label{tab:positioning}
\begin{tabular}{lccccc}
\toprule
Method & Multi-cam & Modality fusion & Decision making & Interpretable & MEVID eval \\
\midrule
TransReID~\cite{he2021transreid}              & $\circ$ &        &        &        &        \\
CLIP-ReID~\cite{li2023clipreid}                & $\circ$ & text aux &      & $\circ$ &        \\
CAL~\cite{gu2022cal} (clothing-change)         &         & clothes-aware &    &        &        \\
OpenGait~\cite{fan2023opengait} family         &         & gait only &      &        &        \\
HuggingGPT~\cite{shen2023hugginggpt}           &         & LLM-tool & \textbullet & \textbullet &  \\
Visual Programming~\cite{gupta2023visprog}     &         & LLM-tool & \textbullet & \textbullet &  \\
Chameleon~\cite{lu2023chameleon}               &         & LLM-tool & \textbullet & \textbullet &  \\
AssistGPT~\cite{gao2023assistgpt}              & $\circ$ & LLM-tool & \textbullet & \textbullet &  \\
\midrule
\textbf{\method{} (ours)} & \textbullet & gait+Re-ID+VLM & \textbullet & \textbullet & \textbullet \\
\bottomrule
\end{tabular}
\end{table*}
